% Person 3 -- Report section: Introduction and Related Work
% IEEEtran conference format. Include with % Person 3 -- Report section: Introduction and Related Work
% IEEEtran conference format. Include with % Person 3 -- Report section: Introduction and Related Work
% IEEEtran conference format. Include with % Person 3 -- Report section: Introduction and Related Work
% IEEEtran conference format. Include with \input{intro_related} and add refs.bib.
% Extractor numbers: review/claim_review_v2.xlsx (4 Nov 2026), scored BEFORE the fixes that
% followed it, so the text says "an earlier version". Claim/planted counts: rebuild of 8 Nov.
% TODO when Person 5 freezes results: add the headline agent-vs-baseline result to the contributions.
% TODO team: contributions 2-4 (tests, agents, dashboard, open release) describe other people's
%      work; check they match what was actually built before submitting.
% Labels stay described as "labelled with AI assistance" (no hand check was done; the 30-row
% spot-check in review/spotcheck.xlsx was a second AI-assisted pass, 30/30 agreed).

\section{Introduction}
\label{sec:intro}

A large share of India's farmland is rainfed, so a late or broken south-west
monsoon can decide a district's harvest. To prepare for such years, ICAR's Central Research
Institute for Dryland Agriculture (ICAR-CRIDA), together with the state agricultural
universities and the Department of Agriculture, has written a District Agriculture
Contingency Plan for almost every district in the country: more than 580 by
2016~\cite{rao2016realtime} and about 650 today~\cite{crida_plans}. Each plan follows
the same template. It records the district's normal rainfall and onset week, its main
crops and sowing windows, and then says what farmers should do when the monsoon is
delayed by two to eight weeks, when a dry spell follows sowing, when drought strikes
early, mid or late in the season, and when floods or heat waves occur.

These plans are static documents built on the data of their time. Of the 105 plans
used in this study, 93 are dated 2011 or 2012, and the crop statistics they cite date
from 2000 to 2011, most often 2008--09. Since then the monsoon itself has been shown to change:
extreme rain events over central India have become more frequent and
stronger~\cite{goswami2006extreme,roxy2017threefold}, dry spells during the peak
monsoon months have become more frequent~\cite{singh2014spells}, and trends differ
from one subdivision to the next~\cite{guhathakurta2008trends}. Cropping patterns have
shifted too. A plan that names a crop the district no longer grows, or a sowing window
that the rain no longer opens, gives poor advice in exactly the year it is needed. The
question is pressing: on 23 June 2026, with about 315 districts flagged for a possibly
weak monsoon, the Union Agriculture Minister asked the states to review and update
their contingency plans for local conditions~\cite{pib2026kharif}. By August, ICAR-CRIDA
had updated the plans of 342 districts in 15 states, but the updated plans are kept at
district level and not maintained centrally~\cite{loksabha2026dacp}, so the plans that can
be studied are still the published versions, nearly all from 2011--2012. Yet we found no study that checks
the plans against newer data at scale.

Doing this by hand is slow. The 105 plans in our study run to 3,321 pages of
semi-structured tables and prose, and every statement in them would have to be
matched to rainfall and crop records. Large language models (LLMs) can read such text,
but used alone they tend to invent numbers and evidence. We therefore split the work.
Code tools compute every number from 45 years of IMD gridded daily
rainfall~\cite{pai2014gridded,imd_gridded} and 26 years of district crop
statistics~\cite{moa_crops}. LLM agents only decide which evidence a claim needs,
judge it, and check each other. Before any agent runs, a deterministic \emph{Claim
Extractor} reads the plans' fixed table template and turns each plan into small, testable
claims, each with a page reference and a word-for-word quote that is checked against the
page. An \emph{Investigator} calls the tools. A \emph{Judge} labels each claim supported,
contradicted or not enough evidence. A \emph{Verifier} re-runs every cited check, and
a \emph{Synthesizer} turns the labels into a 0--100 freshness score for each district.

We ask: \emph{can tool-using LLM agents, which turn plan text into testable claims and
check them against rainfall and crop records, find outdated plan assumptions more
accurately than a rule-based checker or a single LLM prompt?}

Our contributions are:
\begin{itemize}
  \item A dataset of about 16,500 claims from 105 district contingency plans in
        Karnataka, Maharashtra, Rajasthan and Telangana, linked to Census 2011 districts,
        in which every claim carries a page reference and a verbatim quote. On 300
        claims from 30 plans not used during development, labelled with AI assistance
        against the source PDFs, an earlier version of the extractor reached 93.0\%
        precision, and 89.7\% recall on 30 pages of those plans; the errors found there
        were then fixed. The dataset also includes 595 planted errors with known answers
        for testing the checkers.
  \item A set of code tools that compare 1981--2010 with 2011 onward for each
        district: seasonal rainfall, sowing-rain onset, dry spells, heavy-rain days,
        crop-area ranks and crop shifts in deficient years, with non-parametric tests
        corrected for testing many districts.
  \item A multi-agent checker in which no agent computes a number, evaluated
        against a rule-based checker, a single LLM prompt and an ablation without the
        Verifier, using label F1, planted errors caught, wrong numbers per 100 claims,
        citation accuracy and cost per plan.
  \item A per-district freshness score with evidence-cited draft updates for
        experts, shown on a map dashboard, with the data and code released openly.
\end{itemize}

Section~\ref{sec:related} reviews related work. The following sections describe the
data, the evidence tools and the agent system, and then report validation results,
discussion and limitations.

\section{Related Work}
\label{sec:related}

\subsection{Agricultural contingency planning in India}
District contingency plans were prepared by ICAR-CRIDA and its partners, mostly in
2011--2013, and are published as one PDF per district~\cite{crida_plans}. Srinivasa
Rao~\cite{rao2016realtime} describes how they are put into practice: interface
meetings with state departments before each season, and real-time contingency
measures in farmers' fields through the AICRPDA and NICRA networks. ICRISAT's ISSCA
programme lists district contingency and drought-proofing plans as a scalable
adaptation measure~\cite{icrisat_issca}. CRIDA's online contingency tool shows a
district's plan beside current-season rainfall~\cite{crida_plans}, and
SukhaRakshak AI, a chatbot from IWMI, answers farmers' questions from the
plans~\cite{iwmi2025sukha}. All of these use the plans as written. None tests
whether a plan's normal rainfall, crops, sowing windows or trigger thresholds still
match the district's recent data, which is the gap this work addresses.

\subsection{Observed changes in the Indian monsoon}
The IMD $0.25^\circ$ gridded daily rainfall data set~\cite{pai2014gridded} made
long, district-scale rainfall analysis possible, and it is the rainfall source we
use. Studies on IMD data report no long-term trend in all-India monsoon rainfall but
clear changes at subdivision scale~\cite{guhathakurta2008trends}; more frequent and
more intense extreme rain over central
India~\cite{goswami2006extreme,roxy2017threefold}; and more frequent dry spells with
more intense wet spells~\cite{singh2014spells}. This body of work describes how the
monsoon has changed, but it is not tied to agricultural advice. We use the same kind
of indices (onset, dry spells, heavy-rain days, deficient years) to test specific plan
triggers and recommendations, district by district.

\subsection{Automated claim verification}
Fact verification is usually framed as checking a claim against a trusted source.
FEVER~\cite{thorne2018fever} checks claims against Wikipedia sentences,
Climate-FEVER~\cite{diggelmann2020climatefever} does the same for real-world climate
claims, SciFact~\cite{wadden2020scifact} checks scientific claims against paper
abstracts, and TabFact~\cite{chen2020tabfact} checks claims against a single table.
FActScore~\cite{min2023factscore} splits long text into atomic facts and scores each
one, which is close to how we split plans into claims, although our extractor uses rules
over the plans' fixed template rather than an LLM. (On a nine-page sample, an LLM
extractor, Gemini~3.5 Flash, gave the same values wherever both extracted a fact, but
left one page's dates empty and returned only 3 of the 33 rows of two other contingency
tables.) In all of these benchmarks the evidence is
a passage or a small table, and the label is fixed in time. Our claims are checked
against long time series, and the question is whether a statement that was true
for the 2000s still holds since 2011. That requires computation rather than lookup.

\subsection{Tool-using and self-checking LLM agents}
ReAct~\cite{yao2023react} interleaves reasoning steps with tool calls and showed
fewer hallucinations on fact verification when the model could query an external
source. Toolformer~\cite{schick2023toolformer} showed that LLMs can learn when and how
to call APIs such as a calculator. Chain-of-Verification~\cite{dhuliawala2024cove}
reduces hallucination by having the model plan and answer its own verification
questions before revising its answer. In the climate domain,
ChatClimate~\cite{vaghefi2023chatclimate} grounds GPT-4 answers in the IPCC AR6
reports. We build on these ideas but go one step further on numbers: our agents never
compute a value themselves; every number comes from a code tool, and a separate
Verifier re-runs the cited tool calls and sends disagreements back for re-checking.
We found no earlier use of such a system to audit government agricultural plans
against long-term observational data.
 and add refs.bib.
% Extractor numbers: review/claim_review_v2.xlsx (4 Nov 2026), scored BEFORE the fixes that
% followed it, so the text says "an earlier version". Claim/planted counts: rebuild of 8 Nov.
% TODO when Person 5 freezes results: add the headline agent-vs-baseline result to the contributions.
% TODO team: contributions 2-4 (tests, agents, dashboard, open release) describe other people's
%      work; check they match what was actually built before submitting.
% Labels stay described as "labelled with AI assistance" (no hand check was done; the 30-row
% spot-check in review/spotcheck.xlsx was a second AI-assisted pass, 30/30 agreed).

\section{Introduction}
\label{sec:intro}

A large share of India's farmland is rainfed, so a late or broken south-west
monsoon can decide a district's harvest. To prepare for such years, ICAR's Central Research
Institute for Dryland Agriculture (ICAR-CRIDA), together with the state agricultural
universities and the Department of Agriculture, has written a District Agriculture
Contingency Plan for almost every district in the country: more than 580 by
2016~\cite{rao2016realtime} and about 650 today~\cite{crida_plans}. Each plan follows
the same template. It records the district's normal rainfall and onset week, its main
crops and sowing windows, and then says what farmers should do when the monsoon is
delayed by two to eight weeks, when a dry spell follows sowing, when drought strikes
early, mid or late in the season, and when floods or heat waves occur.

These plans are static documents built on the data of their time. Of the 105 plans
used in this study, 93 are dated 2011 or 2012, and the crop statistics they cite date
from 2000 to 2011, most often 2008--09. Since then the monsoon itself has been shown to change:
extreme rain events over central India have become more frequent and
stronger~\cite{goswami2006extreme,roxy2017threefold}, dry spells during the peak
monsoon months have become more frequent~\cite{singh2014spells}, and trends differ
from one subdivision to the next~\cite{guhathakurta2008trends}. Cropping patterns have
shifted too. A plan that names a crop the district no longer grows, or a sowing window
that the rain no longer opens, gives poor advice in exactly the year it is needed. The
question is pressing: on 23 June 2026, with about 315 districts flagged for a possibly
weak monsoon, the Union Agriculture Minister asked the states to review and update
their contingency plans for local conditions~\cite{pib2026kharif}. By August, ICAR-CRIDA
had updated the plans of 342 districts in 15 states, but the updated plans are kept at
district level and not maintained centrally~\cite{loksabha2026dacp}, so the plans that can
be studied are still the published versions, nearly all from 2011--2012. Yet we found no study that checks
the plans against newer data at scale.

Doing this by hand is slow. The 105 plans in our study run to 3,321 pages of
semi-structured tables and prose, and every statement in them would have to be
matched to rainfall and crop records. Large language models (LLMs) can read such text,
but used alone they tend to invent numbers and evidence. We therefore split the work.
Code tools compute every number from 45 years of IMD gridded daily
rainfall~\cite{pai2014gridded,imd_gridded} and 26 years of district crop
statistics~\cite{moa_crops}. LLM agents only decide which evidence a claim needs,
judge it, and check each other. Before any agent runs, a deterministic \emph{Claim
Extractor} reads the plans' fixed table template and turns each plan into small, testable
claims, each with a page reference and a word-for-word quote that is checked against the
page. An \emph{Investigator} calls the tools. A \emph{Judge} labels each claim supported,
contradicted or not enough evidence. A \emph{Verifier} re-runs every cited check, and
a \emph{Synthesizer} turns the labels into a 0--100 freshness score for each district.

We ask: \emph{can tool-using LLM agents, which turn plan text into testable claims and
check them against rainfall and crop records, find outdated plan assumptions more
accurately than a rule-based checker or a single LLM prompt?}

Our contributions are:
\begin{itemize}
  \item A dataset of about 16,500 claims from 105 district contingency plans in
        Karnataka, Maharashtra, Rajasthan and Telangana, linked to Census 2011 districts,
        in which every claim carries a page reference and a verbatim quote. On 300
        claims from 30 plans not used during development, labelled with AI assistance
        against the source PDFs, an earlier version of the extractor reached 93.0\%
        precision, and 89.7\% recall on 30 pages of those plans; the errors found there
        were then fixed. The dataset also includes 595 planted errors with known answers
        for testing the checkers.
  \item A set of code tools that compare 1981--2010 with 2011 onward for each
        district: seasonal rainfall, sowing-rain onset, dry spells, heavy-rain days,
        crop-area ranks and crop shifts in deficient years, with non-parametric tests
        corrected for testing many districts.
  \item A multi-agent checker in which no agent computes a number, evaluated
        against a rule-based checker, a single LLM prompt and an ablation without the
        Verifier, using label F1, planted errors caught, wrong numbers per 100 claims,
        citation accuracy and cost per plan.
  \item A per-district freshness score with evidence-cited draft updates for
        experts, shown on a map dashboard, with the data and code released openly.
\end{itemize}

Section~\ref{sec:related} reviews related work. The following sections describe the
data, the evidence tools and the agent system, and then report validation results,
discussion and limitations.

\section{Related Work}
\label{sec:related}

\subsection{Agricultural contingency planning in India}
District contingency plans were prepared by ICAR-CRIDA and its partners, mostly in
2011--2013, and are published as one PDF per district~\cite{crida_plans}. Srinivasa
Rao~\cite{rao2016realtime} describes how they are put into practice: interface
meetings with state departments before each season, and real-time contingency
measures in farmers' fields through the AICRPDA and NICRA networks. ICRISAT's ISSCA
programme lists district contingency and drought-proofing plans as a scalable
adaptation measure~\cite{icrisat_issca}. CRIDA's online contingency tool shows a
district's plan beside current-season rainfall~\cite{crida_plans}, and
SukhaRakshak AI, a chatbot from IWMI, answers farmers' questions from the
plans~\cite{iwmi2025sukha}. All of these use the plans as written. None tests
whether a plan's normal rainfall, crops, sowing windows or trigger thresholds still
match the district's recent data, which is the gap this work addresses.

\subsection{Observed changes in the Indian monsoon}
The IMD $0.25^\circ$ gridded daily rainfall data set~\cite{pai2014gridded} made
long, district-scale rainfall analysis possible, and it is the rainfall source we
use. Studies on IMD data report no long-term trend in all-India monsoon rainfall but
clear changes at subdivision scale~\cite{guhathakurta2008trends}; more frequent and
more intense extreme rain over central
India~\cite{goswami2006extreme,roxy2017threefold}; and more frequent dry spells with
more intense wet spells~\cite{singh2014spells}. This body of work describes how the
monsoon has changed, but it is not tied to agricultural advice. We use the same kind
of indices (onset, dry spells, heavy-rain days, deficient years) to test specific plan
triggers and recommendations, district by district.

\subsection{Automated claim verification}
Fact verification is usually framed as checking a claim against a trusted source.
FEVER~\cite{thorne2018fever} checks claims against Wikipedia sentences,
Climate-FEVER~\cite{diggelmann2020climatefever} does the same for real-world climate
claims, SciFact~\cite{wadden2020scifact} checks scientific claims against paper
abstracts, and TabFact~\cite{chen2020tabfact} checks claims against a single table.
FActScore~\cite{min2023factscore} splits long text into atomic facts and scores each
one, which is close to how we split plans into claims, although our extractor uses rules
over the plans' fixed template rather than an LLM. (On a nine-page sample, an LLM
extractor, Gemini~3.5 Flash, gave the same values wherever both extracted a fact, but
left one page's dates empty and returned only 3 of the 33 rows of two other contingency
tables.) In all of these benchmarks the evidence is
a passage or a small table, and the label is fixed in time. Our claims are checked
against long time series, and the question is whether a statement that was true
for the 2000s still holds since 2011. That requires computation rather than lookup.

\subsection{Tool-using and self-checking LLM agents}
ReAct~\cite{yao2023react} interleaves reasoning steps with tool calls and showed
fewer hallucinations on fact verification when the model could query an external
source. Toolformer~\cite{schick2023toolformer} showed that LLMs can learn when and how
to call APIs such as a calculator. Chain-of-Verification~\cite{dhuliawala2024cove}
reduces hallucination by having the model plan and answer its own verification
questions before revising its answer. In the climate domain,
ChatClimate~\cite{vaghefi2023chatclimate} grounds GPT-4 answers in the IPCC AR6
reports. We build on these ideas but go one step further on numbers: our agents never
compute a value themselves; every number comes from a code tool, and a separate
Verifier re-runs the cited tool calls and sends disagreements back for re-checking.
We found no earlier use of such a system to audit government agricultural plans
against long-term observational data.
 and add refs.bib.
% Extractor numbers: review/claim_review_v2.xlsx (4 Nov 2026), scored BEFORE the fixes that
% followed it, so the text says "an earlier version". Claim/planted counts: rebuild of 8 Nov.
% TODO when Person 5 freezes results: add the headline agent-vs-baseline result to the contributions.
% TODO team: contributions 2-4 (tests, agents, dashboard, open release) describe other people's
%      work; check they match what was actually built before submitting.
% Labels stay described as "labelled with AI assistance" (no hand check was done; the 30-row
% spot-check in review/spotcheck.xlsx was a second AI-assisted pass, 30/30 agreed).

\section{Introduction}
\label{sec:intro}

A large share of India's farmland is rainfed, so a late or broken south-west
monsoon can decide a district's harvest. To prepare for such years, ICAR's Central Research
Institute for Dryland Agriculture (ICAR-CRIDA), together with the state agricultural
universities and the Department of Agriculture, has written a District Agriculture
Contingency Plan for almost every district in the country: more than 580 by
2016~\cite{rao2016realtime} and about 650 today~\cite{crida_plans}. Each plan follows
the same template. It records the district's normal rainfall and onset week, its main
crops and sowing windows, and then says what farmers should do when the monsoon is
delayed by two to eight weeks, when a dry spell follows sowing, when drought strikes
early, mid or late in the season, and when floods or heat waves occur.

These plans are static documents built on the data of their time. Of the 105 plans
used in this study, 93 are dated 2011 or 2012, and the crop statistics they cite date
from 2000 to 2011, most often 2008--09. Since then the monsoon itself has been shown to change:
extreme rain events over central India have become more frequent and
stronger~\cite{goswami2006extreme,roxy2017threefold}, dry spells during the peak
monsoon months have become more frequent~\cite{singh2014spells}, and trends differ
from one subdivision to the next~\cite{guhathakurta2008trends}. Cropping patterns have
shifted too. A plan that names a crop the district no longer grows, or a sowing window
that the rain no longer opens, gives poor advice in exactly the year it is needed. The
question is pressing: on 23 June 2026, with about 315 districts flagged for a possibly
weak monsoon, the Union Agriculture Minister asked the states to review and update
their contingency plans for local conditions~\cite{pib2026kharif}. By August, ICAR-CRIDA
had updated the plans of 342 districts in 15 states, but the updated plans are kept at
district level and not maintained centrally~\cite{loksabha2026dacp}, so the plans that can
be studied are still the published versions, nearly all from 2011--2012. Yet we found no study that checks
the plans against newer data at scale.

Doing this by hand is slow. The 105 plans in our study run to 3,321 pages of
semi-structured tables and prose, and every statement in them would have to be
matched to rainfall and crop records. Large language models (LLMs) can read such text,
but used alone they tend to invent numbers and evidence. We therefore split the work.
Code tools compute every number from 45 years of IMD gridded daily
rainfall~\cite{pai2014gridded,imd_gridded} and 26 years of district crop
statistics~\cite{moa_crops}. LLM agents only decide which evidence a claim needs,
judge it, and check each other. Before any agent runs, a deterministic \emph{Claim
Extractor} reads the plans' fixed table template and turns each plan into small, testable
claims, each with a page reference and a word-for-word quote that is checked against the
page. An \emph{Investigator} calls the tools. A \emph{Judge} labels each claim supported,
contradicted or not enough evidence. A \emph{Verifier} re-runs every cited check, and
a \emph{Synthesizer} turns the labels into a 0--100 freshness score for each district.

We ask: \emph{can tool-using LLM agents, which turn plan text into testable claims and
check them against rainfall and crop records, find outdated plan assumptions more
accurately than a rule-based checker or a single LLM prompt?}

Our contributions are:
\begin{itemize}
  \item A dataset of about 16,500 claims from 105 district contingency plans in
        Karnataka, Maharashtra, Rajasthan and Telangana, linked to Census 2011 districts,
        in which every claim carries a page reference and a verbatim quote. On 300
        claims from 30 plans not used during development, labelled with AI assistance
        against the source PDFs, an earlier version of the extractor reached 93.0\%
        precision, and 89.7\% recall on 30 pages of those plans; the errors found there
        were then fixed. The dataset also includes 595 planted errors with known answers
        for testing the checkers.
  \item A set of code tools that compare 1981--2010 with 2011 onward for each
        district: seasonal rainfall, sowing-rain onset, dry spells, heavy-rain days,
        crop-area ranks and crop shifts in deficient years, with non-parametric tests
        corrected for testing many districts.
  \item A multi-agent checker in which no agent computes a number, evaluated
        against a rule-based checker, a single LLM prompt and an ablation without the
        Verifier, using label F1, planted errors caught, wrong numbers per 100 claims,
        citation accuracy and cost per plan.
  \item A per-district freshness score with evidence-cited draft updates for
        experts, shown on a map dashboard, with the data and code released openly.
\end{itemize}

Section~\ref{sec:related} reviews related work. The following sections describe the
data, the evidence tools and the agent system, and then report validation results,
discussion and limitations.

\section{Related Work}
\label{sec:related}

\subsection{Agricultural contingency planning in India}
District contingency plans were prepared by ICAR-CRIDA and its partners, mostly in
2011--2013, and are published as one PDF per district~\cite{crida_plans}. Srinivasa
Rao~\cite{rao2016realtime} describes how they are put into practice: interface
meetings with state departments before each season, and real-time contingency
measures in farmers' fields through the AICRPDA and NICRA networks. ICRISAT's ISSCA
programme lists district contingency and drought-proofing plans as a scalable
adaptation measure~\cite{icrisat_issca}. CRIDA's online contingency tool shows a
district's plan beside current-season rainfall~\cite{crida_plans}, and
SukhaRakshak AI, a chatbot from IWMI, answers farmers' questions from the
plans~\cite{iwmi2025sukha}. All of these use the plans as written. None tests
whether a plan's normal rainfall, crops, sowing windows or trigger thresholds still
match the district's recent data, which is the gap this work addresses.

\subsection{Observed changes in the Indian monsoon}
The IMD $0.25^\circ$ gridded daily rainfall data set~\cite{pai2014gridded} made
long, district-scale rainfall analysis possible, and it is the rainfall source we
use. Studies on IMD data report no long-term trend in all-India monsoon rainfall but
clear changes at subdivision scale~\cite{guhathakurta2008trends}; more frequent and
more intense extreme rain over central
India~\cite{goswami2006extreme,roxy2017threefold}; and more frequent dry spells with
more intense wet spells~\cite{singh2014spells}. This body of work describes how the
monsoon has changed, but it is not tied to agricultural advice. We use the same kind
of indices (onset, dry spells, heavy-rain days, deficient years) to test specific plan
triggers and recommendations, district by district.

\subsection{Automated claim verification}
Fact verification is usually framed as checking a claim against a trusted source.
FEVER~\cite{thorne2018fever} checks claims against Wikipedia sentences,
Climate-FEVER~\cite{diggelmann2020climatefever} does the same for real-world climate
claims, SciFact~\cite{wadden2020scifact} checks scientific claims against paper
abstracts, and TabFact~\cite{chen2020tabfact} checks claims against a single table.
FActScore~\cite{min2023factscore} splits long text into atomic facts and scores each
one, which is close to how we split plans into claims, although our extractor uses rules
over the plans' fixed template rather than an LLM. (On a nine-page sample, an LLM
extractor, Gemini~3.5 Flash, gave the same values wherever both extracted a fact, but
left one page's dates empty and returned only 3 of the 33 rows of two other contingency
tables.) In all of these benchmarks the evidence is
a passage or a small table, and the label is fixed in time. Our claims are checked
against long time series, and the question is whether a statement that was true
for the 2000s still holds since 2011. That requires computation rather than lookup.

\subsection{Tool-using and self-checking LLM agents}
ReAct~\cite{yao2023react} interleaves reasoning steps with tool calls and showed
fewer hallucinations on fact verification when the model could query an external
source. Toolformer~\cite{schick2023toolformer} showed that LLMs can learn when and how
to call APIs such as a calculator. Chain-of-Verification~\cite{dhuliawala2024cove}
reduces hallucination by having the model plan and answer its own verification
questions before revising its answer. In the climate domain,
ChatClimate~\cite{vaghefi2023chatclimate} grounds GPT-4 answers in the IPCC AR6
reports. We build on these ideas but go one step further on numbers: our agents never
compute a value themselves; every number comes from a code tool, and a separate
Verifier re-runs the cited tool calls and sends disagreements back for re-checking.
We found no earlier use of such a system to audit government agricultural plans
against long-term observational data.
 and add refs.bib.
% Extractor numbers: review/claim_review_v2.xlsx (4 Nov 2026), scored BEFORE the fixes that
% followed it, so the text says "an earlier version". Claim/planted counts: rebuild of 8 Nov.
% TODO when Person 5 freezes results: add the headline agent-vs-baseline result to the contributions.
% TODO team: contributions 2-4 (tests, agents, dashboard, open release) describe other people's
%      work; check they match what was actually built before submitting.
% Labels stay described as "labelled with AI assistance" (no hand check was done; the 30-row
% spot-check in review/spotcheck.xlsx was a second AI-assisted pass, 30/30 agreed).

\section{Introduction}
\label{sec:intro}

A large share of India's farmland is rainfed, so a late or broken south-west
monsoon can decide a district's harvest. To prepare for such years, ICAR's Central Research
Institute for Dryland Agriculture (ICAR-CRIDA), together with the state agricultural
universities and the Department of Agriculture, has written a District Agriculture
Contingency Plan for almost every district in the country: more than 580 by
2016~\cite{rao2016realtime} and about 650 today~\cite{crida_plans}. Each plan follows
the same template. It records the district's normal rainfall and onset week, its main
crops and sowing windows, and then says what farmers should do when the monsoon is
delayed by two to eight weeks, when a dry spell follows sowing, when drought strikes
early, mid or late in the season, and when floods or heat waves occur.

These plans are static documents built on the data of their time. Of the 105 plans
used in this study, 93 are dated 2011 or 2012, and the crop statistics they cite date
from 2000 to 2011, most often 2008--09. Since then the monsoon itself has been shown to change:
extreme rain events over central India have become more frequent and
stronger~\cite{goswami2006extreme,roxy2017threefold}, dry spells during the peak
monsoon months have become more frequent~\cite{singh2014spells}, and trends differ
from one subdivision to the next~\cite{guhathakurta2008trends}. Cropping patterns have
shifted too. A plan that names a crop the district no longer grows, or a sowing window
that the rain no longer opens, gives poor advice in exactly the year it is needed. The
question is pressing: on 23 June 2026, with about 315 districts flagged for a possibly
weak monsoon, the Union Agriculture Minister asked the states to review and update
their contingency plans for local conditions~\cite{pib2026kharif}. By August, ICAR-CRIDA
had updated the plans of 342 districts in 15 states, but the updated plans are kept at
district level and not maintained centrally~\cite{loksabha2026dacp}, so the plans that can
be studied are still the published versions, nearly all from 2011--2012. Yet we found no study that checks
the plans against newer data at scale.

Doing this by hand is slow. The 105 plans in our study run to 3,321 pages of
semi-structured tables and prose, and every statement in them would have to be
matched to rainfall and crop records. Large language models (LLMs) can read such text,
but used alone they tend to invent numbers and evidence. We therefore split the work.
Code tools compute every number from 45 years of IMD gridded daily
rainfall~\cite{pai2014gridded,imd_gridded} and 26 years of district crop
statistics~\cite{moa_crops}. LLM agents only decide which evidence a claim needs,
judge it, and check each other. Before any agent runs, a deterministic \emph{Claim
Extractor} reads the plans' fixed table template and turns each plan into small, testable
claims, each with a page reference and a word-for-word quote that is checked against the
page. An \emph{Investigator} calls the tools. A \emph{Judge} labels each claim supported,
contradicted or not enough evidence. A \emph{Verifier} re-runs every cited check, and
a \emph{Synthesizer} turns the labels into a 0--100 freshness score for each district.

We ask: \emph{can tool-using LLM agents, which turn plan text into testable claims and
check them against rainfall and crop records, find outdated plan assumptions more
accurately than a rule-based checker or a single LLM prompt?}

Our contributions are:
\begin{itemize}
  \item A dataset of about 16,500 claims from 105 district contingency plans in
        Karnataka, Maharashtra, Rajasthan and Telangana, linked to Census 2011 districts,
        in which every claim carries a page reference and a verbatim quote. On 300
        claims from 30 plans not used during development, labelled with AI assistance
        against the source PDFs, an earlier version of the extractor reached 93.0\%
        precision, and 89.7\% recall on 30 pages of those plans; the errors found there
        were then fixed. The dataset also includes 595 planted errors with known answers
        for testing the checkers.
  \item A set of code tools that compare 1981--2010 with 2011 onward for each
        district: seasonal rainfall, sowing-rain onset, dry spells, heavy-rain days,
        crop-area ranks and crop shifts in deficient years, with non-parametric tests
        corrected for testing many districts.
  \item A multi-agent checker in which no agent computes a number, evaluated
        against a rule-based checker, a single LLM prompt and an ablation without the
        Verifier, using label F1, planted errors caught, wrong numbers per 100 claims,
        citation accuracy and cost per plan.
  \item A per-district freshness score with evidence-cited draft updates for
        experts, shown on a map dashboard, with the data and code released openly.
\end{itemize}

Section~\ref{sec:related} reviews related work. The following sections describe the
data, the evidence tools and the agent system, and then report validation results,
discussion and limitations.

\section{Related Work}
\label{sec:related}

\subsection{Agricultural contingency planning in India}
District contingency plans were prepared by ICAR-CRIDA and its partners, mostly in
2011--2013, and are published as one PDF per district~\cite{crida_plans}. Srinivasa
Rao~\cite{rao2016realtime} describes how they are put into practice: interface
meetings with state departments before each season, and real-time contingency
measures in farmers' fields through the AICRPDA and NICRA networks. ICRISAT's ISSCA
programme lists district contingency and drought-proofing plans as a scalable
adaptation measure~\cite{icrisat_issca}. CRIDA's online contingency tool shows a
district's plan beside current-season rainfall~\cite{crida_plans}, and
SukhaRakshak AI, a chatbot from IWMI, answers farmers' questions from the
plans~\cite{iwmi2025sukha}. All of these use the plans as written. None tests
whether a plan's normal rainfall, crops, sowing windows or trigger thresholds still
match the district's recent data, which is the gap this work addresses.

\subsection{Observed changes in the Indian monsoon}
The IMD $0.25^\circ$ gridded daily rainfall data set~\cite{pai2014gridded} made
long, district-scale rainfall analysis possible, and it is the rainfall source we
use. Studies on IMD data report no long-term trend in all-India monsoon rainfall but
clear changes at subdivision scale~\cite{guhathakurta2008trends}; more frequent and
more intense extreme rain over central
India~\cite{goswami2006extreme,roxy2017threefold}; and more frequent dry spells with
more intense wet spells~\cite{singh2014spells}. This body of work describes how the
monsoon has changed, but it is not tied to agricultural advice. We use the same kind
of indices (onset, dry spells, heavy-rain days, deficient years) to test specific plan
triggers and recommendations, district by district.

\subsection{Automated claim verification}
Fact verification is usually framed as checking a claim against a trusted source.
FEVER~\cite{thorne2018fever} checks claims against Wikipedia sentences,
Climate-FEVER~\cite{diggelmann2020climatefever} does the same for real-world climate
claims, SciFact~\cite{wadden2020scifact} checks scientific claims against paper
abstracts, and TabFact~\cite{chen2020tabfact} checks claims against a single table.
FActScore~\cite{min2023factscore} splits long text into atomic facts and scores each
one, which is close to how we split plans into claims, although our extractor uses rules
over the plans' fixed template rather than an LLM. (On a nine-page sample, an LLM
extractor, Gemini~3.5 Flash, gave the same values wherever both extracted a fact, but
left one page's dates empty and returned only 3 of the 33 rows of two other contingency
tables.) In all of these benchmarks the evidence is
a passage or a small table, and the label is fixed in time. Our claims are checked
against long time series, and the question is whether a statement that was true
for the 2000s still holds since 2011. That requires computation rather than lookup.

\subsection{Tool-using and self-checking LLM agents}
ReAct~\cite{yao2023react} interleaves reasoning steps with tool calls and showed
fewer hallucinations on fact verification when the model could query an external
source. Toolformer~\cite{schick2023toolformer} showed that LLMs can learn when and how
to call APIs such as a calculator. Chain-of-Verification~\cite{dhuliawala2024cove}
reduces hallucination by having the model plan and answer its own verification
questions before revising its answer. In the climate domain,
ChatClimate~\cite{vaghefi2023chatclimate} grounds GPT-4 answers in the IPCC AR6
reports. We build on these ideas but go one step further on numbers: our agents never
compute a value themselves; every number comes from a code tool, and a separate
Verifier re-runs the cited tool calls and sends disagreements back for re-checking.
We found no earlier use of such a system to audit government agricultural plans
against long-term observational data.
